\chapter{Muestreo en tiempo}

\section{Introducci\'{o}n}

\subsection{Definiciones b\'{a}sicas}

\begin{enumerate}
\item   Recordemos la definici\'{o}n \textbf{de la sucesi\'{o}n de funciones} cuyo l\'{\i}mite es el
impulso unitario
\begin{equation}
  r_{\Delta} \left( t - t_{0} \right) = 
  \left \{ \begin{array}{ll}
       0                       &       t < t_{0},                                              \\
  \\
       \frac{1}{\Delta}        &       t_{0} \leq      t       \leq    t_{0} + \Delta,         \\
  \\
       0                       &       t_{0} + \Delta < t 
  \end{array} \right.
\end{equation}
                donde $\Delta$ es una constante real no negativa.

\item   Supongamos que tenemos una se\~{n}al $x \left( t \right)$. Podemos definir una nueva se\~{n}al
\begin{equation}
                        z \left( t \right) = x \left( t \right) \delta \left( t - t_{0}\right).
                \end{equation}
                
\item Para obtener los valores de la funci\'{o}n $z(t)$, construyamos una sucesi\'{o}n de funciones
                \begin{equation}
                        z_{\Delta} \left( t \right) = \left \{
                        \begin{array}{ll}
                 0                                       &       t < t_{0},                              \\
                 z\left( t \right) \frac{1}{\Delta}      &       t_{0} \leq t \leq t_{0} + \Delta,       \\
                 0                                       &       t_{0} + \Delta < t.
                        \end{array}
                        \right.
                \end{equation}
                La funci\'{o}n $z \left( t \right)$ es el l\'{\i}mite de la sucesi\'{o}n cuando
                $\Delta \rightarrow 0$, con la siguiente definici\'{o}n
                \begin{equation}
                        z\left(t \right) = \left \{
                        \begin{array}{ll}
                                0                       &       t \neq t_{0},   \\
                                x\left(t_{0}\right)     &       t  = t_{0}.
                        \end{array}             \label{muescont}
                        \right.
                \end{equation}

\item   Podemos resumir la operaci\'{o}n realizada en la ec. (\ref{muescont}) como              
                \begin{equation}
                        z \left( t \right) = 
                                x\left( t\right) \delta \left( t-t_{0}\right) =
                                x\left( t_{0}\right) \delta \left( t-t_{0}\right) .
                \end{equation}

\end{enumerate}

\section{Muestreo por modulaci\'{o}n con un tren de impulsos, tiempo continuo}

\begin{enumerate}
\item Veamos la Fig. (\ref{samponeeps}), y recordemos la definici\'{o}n de un tren de impulsos de tiempo 
continuo
\begin{equation}
        p\left( t\right) =\sum_{n=-\infty }^{\infty }\delta \left( t-nT\right) .        \label{samptemp}
\end{equation}
Sabemos adem\'{a}s que la transformada de Fourier de $p(t)$ es
\begin{equation}
        P\left( j\omega \right) =
                \frac{2\pi }{T}\sum_{k=-\infty }^{\infty }\delta\left( \omega -k\frac{2\pi }{T}\right) .
                \label{fouriertrenimp}
\end{equation}

\item Calculemos ahora el producto de $p\left( t\right) $ con una se\~{n}al $x\left( t\right) $
\begin{equation}
        x_{p}\left( t\right) = x\left( t\right) p\left( t\right) = 
                  \sum_{n=-\infty }^{\infty } x(t) \delta \left( t-nT\right).
\end{equation}
Por la propiedad del producto de una se\~{n}al con un impulso sabemos que
\begin{equation}
        x\left( t\right) \delta \left( t-t_{0}\right) =x\left( t_{0}\right) \delta \left( t-t_{0}\right) .
\end{equation}
Por lo tanto resulta 
\begin{equation}
        x_{p}\left( t\right) =\sum_{n=-\infty }^{\infty }x\left( nT\right) \delta \left( t-nT\right) .
\end{equation}

\item Por la propiedad de la multiplicaci\'{o}n de la Transformada de Fourier
sabemos que la transformada de Fourier de la se\~{n}al $x_{p}(t)$ debe ser
\begin{equation}
        X_{p}\left( j\omega \right) =
                \frac{1}{2\pi }
                \int_{-\infty }^{\infty}
                        X\left( j\theta \right) P\left( j\left( \omega -\theta \right) \right) d\theta .
\end{equation}
Teniendo en cuenta la ecuaci\'{o}n (\ref{fouriertrenimp}) escribimos
\begin{equation}
        X_{p}\left( j\omega \right) =
\frac{1}{2\pi }
\int_{-\infty }^{\infty}
  \frac{2\pi }{T}
  \sum_{k=-\infty }^{\infty }    
      X\left( j\theta \right) \delta\left( \omega -k\frac{2\pi}{T} -\theta \right)  
d\theta,
\end{equation}
donde, intercambiando las posiciones de la integral y la sumatoria y simplificando, finalmente 
permite deducir
\begin{equation}
        X_{p}\left( j\omega \right) =
  \frac{1}{T}
  \sum_{k=-\infty }^{\infty }
  \int_{-\infty }^{\infty}
      X\left( j\theta \right) \delta\left( \omega -k\frac{2\pi}{T} -\theta \right)  
  d\theta.
\end{equation}

\item Dada la propiedad de la convoluci\'{o}n de una se\~{n}al con un impulso
\begin{equation}
 X \left( j\omega \right) * \delta \left( \omega -k \frac{2\pi}{T} \right) =
  X\left( j\omega - k \frac{2 \pi}{T} \right),
\end{equation}
resulta 
\begin{equation}
        X_{p}\left( j\omega \right) =
 \frac{1}{T} 
 \sum_{k=-\infty }^{\infty } X\left(j\left( \omega -k\frac{2\pi }{T}\right) \right) .    \label{sampomega}
\end{equation}

\item Para explicar en palabras la ec. (\ref{sampomega}), recordemos que se le da el nombre de 
  espectro de la se\~{n}al $x(t)$ a su transformada de Fourier $X(j \omega)$. 
  La ec. (\ref{sampomega}) es la transformada de Fourier de la se\~{n}al muestreada y 
  es una sumatoria o superposici\'{o}n de copias del espectro original desplazadas en frecuencia a 
  intervalos regulares $\frac{2 \pi}{T}$. 
\end{enumerate}

\begin{figure}
        \centering
        \includegraphics[height=3.5in]{./graphics/sampone.eps}
        \caption{Muestreo en el dominio del tiempo: La se\~{n}al original se modula con 
        tren de impulsos.}
        \label{samponeeps}
\end{figure}

\section{Teorema del muestreo}

\begin{enumerate}
        \item   Sea $x\left( t\right) $ una se\~{n}al cuyo espectro tiene un ancho de banda limitado con 
          $X\left( j\omega \right) =0$ para $\left| \omega \right| >\omega_{a}$, donde $\omega_{a}$ 
          es el ancho de banda. Entonces $x\left( t\right) $ queda correctamente especificada por 
          sus muestras $x\left( nT\right)$, $n=0, 1, 2,....,$ si
        \begin{equation}
                \omega_{s}=\frac{2\pi }{T}>2\omega_{a}.
        \end{equation}
        
        \item En otras palabras, si generamos una se\~{n}al $x_{T}^{ }\left( t\right) $ modulando un tren 
          de impulsos peri\'{o}dicos de frecuencia $\omega_{s}$ con una se\~{n}al $x\left( t\right),$ 
          podemos recuperar la se\~{n}al $x(t)$ haciendo pasar la se\~{n}al $x_{T}\left( t\right)$ por un 
          filtro pasabajos ideal con ganancia $T$ y frecuencia de corte superior a $\omega_{a}$ e inferior 
          a $\omega_{s}-\omega_{a}$.

        \item Las dos ecuaciones que debemos observar para entender la clave del muestreo son 
          (\ref{samptemp}) y (\ref{sampomega}), repetidas a continuaci\'{o}n
        \begin{equation}
                p\left( t\right) =\sum_{n=-\infty }^{\infty }\delta \left( t-nT\right),
        \end{equation}
        \begin{equation}
                X_{p}\left( j\omega \right) =
                \frac{1}{T}
                \sum_{k=-\infty }^{\infty } X\left(j\left( \omega -k\frac{2\pi }{T}\right) \right).
        \end{equation}
        La separaci\'{o}n en el tren de impulsos temporal es $T$, y en el tren de impulsos en frecuencia 
        es $\frac{2\pi }{T}$. Cuando mayor es el tiempo de muestreo, menor es la separaci\'{o}n de los
        centros de las copias del espectro original en frecuencia.

        \item Si los centros de los espectros se acercan demasiado, entonces los extremos de las copias
          del espectro original se superponen y confunden, provocando el fen\'{o}meno de aliasing.

\end{enumerate}
        
\begin{figure}
        \centering
        \includegraphics[height=3.9in]{./graphics/samptwo.eps}
        \caption{Muestreo en el dominio de la frecuencia. La multiplicaci\'{o}n en el dominio del tiempo
        corresponde a una convoluci\'{o}n en el dominio de la frecuencia. La  convoluci\'{o}n de un impulso
        ubicado en $\frac{k 2 \Pi}{T}$ con la se\~{n}al original genera una copia de esta cuya frecuencia
        central es $\frac{k 2 \Pi}{T}$. Obtenemos as\'{\i} infinitas copias del espectro de la se\~{n}al
        original ubicadas en las frecuencias $k \omega_{s}$.}
        \label{samptwoeps}
\end{figure}

\begin{figure}
        \centering
        \includegraphics[height=4.0in]{./graphics/sampthree.eps}
        \caption{Aliasing o enmascaramiento de frecuencias por muestreo demasiado lento.}
        \label{sampthreeeps}
\end{figure}

\section{Reconstrucci\'{o}n ideal de se\~{n}ales muestreadas}

\begin{enumerate}

\item Recordemos la propiedad de la multiplicaci\'{o}n de la Transformada de Fourier
        \begin{equation}
                X(j \omega) \, H(j \omega)) 
                \begin{array}{c} TF \\ \longleftrightarrow \end{array} x(t) * h(t).
        \end{equation}
        
\item Supongamos un  filtro  pasabajo  ideal  cuya respuesta en frecuencia es
        \begin{equation}
                H(j \omega) =
                        \left[  \begin{array}{cc}
                                1       &       |\omega| \leq \omega_{c},       \\
                                0       &       |\omega| > \omega_{c}.
                        \end{array}\right.
        \end{equation}
        Por lo tanto su respuesta al impulso es
        \begin{equation}
                h(t) =   \omega_{c} \, T \, \frac{sin(\omega_{c} \, t)}{\pi \, \omega_{c} \, t}.
        \end{equation}

\item La se\~{n}al  reconstruida es  entonces
        \begin{equation}
                x_{r}\left(  t\right) =
                \sum_{\infty}^{\infty}
                        x(n T) \, \frac{\omega_{c} \, T}{\pi} \frac{sin(\omega_{c} \, (t-nT))}{\omega_{c} \, (t-nT)}.
        \end{equation}

\item Esta es una reconstrucci\'{o}n ideal porque estamos usando un filtro pasabajos ideal, cuya
  respuesta es no causal.

\end{enumerate}

\begin{figure}
  \centering
  \includegraphics[height=3in]{./graphics/muestreocontinua.eps}
  \caption{Una se\~{n}al continua que ser\'{a} muestreada}
  \label{muestreocontinuaeps}
\end{figure}

\begin{figure}
  \centering
  \includegraphics[height=3in]{./graphics/muestreomuestreada.eps}
  \caption{La se\~{n}al muestreada}
  \label{muestreomuestreadaeps}
\end{figure}

\begin{figure}
  \centering
  \includegraphics[height=3in]{./graphics/muestreoreconstruccion.eps}
  \caption{Reconstrucci\'{o}n ideal de la se\~{n}al muestreada}
  \label{muestreoreconstruccioneps}
\end{figure}

\section{Reconstrucci\'{o}n real de se\~{n}ales muestreadas}

Una se\~{n}al continua muestreada se puede reconstruir mediante
pulsos (vea la Figura \ref{muestreoreconstruccionrealeps})

\begin{figure}
  \centering
  \includegraphics[height=3in]{./graphics/contconv222.eps}
  \caption{Reconstrucci\'{o}n real de la se\~{n}al muestreada,
    mediante sumatoria de pulsos}
  \label{muestreoreconstruccionrealeps}
\end{figure}
